% Options for packages loaded elsewhere
\PassOptionsToPackage{unicode}{hyperref}
\PassOptionsToPackage{hyphens}{url}
\PassOptionsToPackage{dvipsnames,svgnames,x11names}{xcolor}
\documentclass[
  10pt,
  fontset=fandol]{ctexart}
\usepackage{xcolor}
\usepackage[margin=16mm]{geometry}
\usepackage{amsmath,amssymb}
\setcounter{secnumdepth}{-\maxdimen} % remove section numbering
\usepackage{iftex}
\ifPDFTeX
  \usepackage[T1]{fontenc}
  \usepackage[utf8]{inputenc}
  \usepackage{textcomp} % provide euro and other symbols
\else % if luatex or xetex
  \usepackage{unicode-math} % this also loads fontspec
  \defaultfontfeatures{Scale=MatchLowercase}
  \defaultfontfeatures[\rmfamily]{Ligatures=TeX,Scale=1}
\fi
\usepackage{lmodern}
\ifPDFTeX\else
  % xetex/luatex font selection
\fi
% Use upquote if available, for straight quotes in verbatim environments
\IfFileExists{upquote.sty}{\usepackage{upquote}}{}
\IfFileExists{microtype.sty}{% use microtype if available
  \usepackage[]{microtype}
  \UseMicrotypeSet[protrusion]{basicmath} % disable protrusion for tt fonts
}{}
\makeatletter
\@ifundefined{KOMAClassName}{% if non-KOMA class
  \IfFileExists{parskip.sty}{%
    \usepackage{parskip}
  }{% else
    \setlength{\parindent}{0pt}
    \setlength{\parskip}{6pt plus 2pt minus 1pt}}
}{% if KOMA class
  \KOMAoptions{parskip=half}}
\makeatother
\setlength{\emergencystretch}{3em} % prevent overfull lines
\providecommand{\tightlist}{%
  \setlength{\itemsep}{0pt}\setlength{\parskip}{0pt}}
\usepackage{amsmath,amssymb,mathtools,needspace}
\xeCJKDeclareCharClass{CJK}{"2032,"0400->"04FF,"2160->"216B,"2460->"2473,"25A0,"25C7,"25CB,"25CF}
\IfFontExistsTF{Arial Unicode MS}{\xeCJKsetup{AutoFallBack=true}\setCJKfallbackfamilyfont{\CJKrmdefault}{Arial Unicode MS}}{}
\setcounter{secnumdepth}{0}
\setlength{\emergencystretch}{2em}
\allowdisplaybreaks
\usepackage{bookmark}
\IfFileExists{xurl.sty}{\usepackage{xurl}}{} % add URL line breaks if available
\urlstyle{same}
\hypersetup{
  pdftitle={小结},
  colorlinks=true,
  linkcolor={Maroon},
  filecolor={Maroon},
  citecolor={Blue},
  urlcolor={Blue},
  pdfcreator={LaTeX via pandoc}}

\title{小结}
\author{}
\date{}

\begin{document}
\maketitle

\subsection{小结}\label{ux5c0fux7ed3}

本章阐述快速 Walsh 变换 FWT 的设计机理与设计方法。可以看到，FWT 本质上是一类二分法，其设计思想是，逐步二分所给计算模型 \(N\)-WT，令其规模 \(N\) 逐次减半，直到规模为 1 时，所归结出的 1-WT 即为所要的结果：

\[
\underset{\text{（计算模型）}}{N\text{-WT}}\Rightarrow2\text{ 个 }N/2\text{-WT}\Rightarrow4\text{ 个 }N/4\text{-WT}\Rightarrow\cdots\Rightarrow\underset{\text{（所求结果）}}{N\text{ 个 }1\text{-WT}}.
\]

注意到 \(N\)-WT 的变换矩阵是 Hadamard 阵 \(H_N\)，上述 FWT 的设计过程本质上是 Hadamard 阵的加工过程

\[
H_N\Rightarrow H_{N/2}\Rightarrow H_{N/4}\Rightarrow\cdots\Rightarrow H_1.
\]

再对比 Hadamard 阵的生成过程（10.3.4 小节）

\[
H_1\Rightarrow H_2\Rightarrow H_4\Rightarrow\cdots\Rightarrow H_N,
\]

可以看到，快速 Walsh 变换的演化过程同 Hadamard 阵的生成过程，它们两者互为反过程。如果后者视为\textbf{进化过程}（矩阵阶数逐步倍增），那么前者则是\textbf{退化过程}（矩阵阶数逐次减半）。在适当定义``规模''的前提下，它们两者全都从属于二分演化模式。

Walsh 分析处处都渗透了对立统一的辩证思维。

正因为 Walsh 函数具有极度的数学美，正由于 Walsh 分析展现了一种新的思维方式，因而在 Walsh 分析的基础上可以开展许多重要的研究。

快速 Walsh 变换是快速变换的一个重要的组成部分。运用变异技术，基于 Walsh 变换可以派生出其他种种快速变换，诸如 Haar 变换、斜变换、Hartly 变换等等，从而实现快速变换方法的大统一。

Walsh 分析有着广泛的应用前景。然而更为重要的是，它展现了一种新的数学方法------\textbf{演化数学方法}。

宇宙是演化的，生物是演化的。时至今日，辩证法关于发展变化的观点，即事物从低级到高级不断演化的观点，已经被科学界认为是无须论证的常识了。

Walsh 函数的演化分析用数学语言表述了这种``常识''。

Walsh 函数的演化分析无疑是新的数学革命即将爆发的先兆。还是 Harmuth 有远见：\textbf{Walsh 分析的研究将导致一场数学革命，就像 17、18 世纪的微积分那样。}

\end{document}
